% from ega4-10-fr.tex
\oldpage[IV-3]{110}
à un point, soit de dimension $1$, et pour tout point fermé $s\in S'$,
$\mathcal O_{S',s}$ est nécessairement de dimension $1$.

On déduit de ces remarques et de (10.6.1) que tout préschéma localement de type
fini sur $S$ vérifie aussi les propriétés $1^\circ$, $2^\circ$ et $3^\circ$ de
(10.6.1)~: il en est ainsi en particulier des préschémas \emph{localement de type
fini sur un corps ou sur $\mathbf Z$}.
\end{env}

\begin{env}[10.7.2]
\label{IV.10.7.2-fr}
Soit $A$ un anneau local noethérien \emph{universellement caténaire} et soit $S$
le \emph{complémentaire dans $X=\operatorname{Spec}(A)$ du point fermé} $a$.
Alors $S$ vérifie les conditions $1^\circ$, $2^\circ$ et $3^\circ$ de (10.6.1)~:
en effet, on a déjà vu que $S$ est un préschéma de Jacobson (10.5.9)~; comme $A$
est universellement caténaire, il en est de même des anneaux locaux $A_{\mathfrak p}$
aux idéaux premiers de $A$ (5.6.3). D'autre part, une composante irréductible $S'$
de $S$ est le complémentaire de $a$ dans une composante irréductible $X'$ de $X$~;
pour tout point fermé $x$ de $S$, l'adhérence de $x$ dans $X$ est donc
$\{x,a\}$, autrement dit $\dim(\overline{\{x\}})=1$ et par suite, puisque $X'$
est par hypothèse biéquidimensionnel ($\mathbf 0$, 14.3.3), on a
$\dim(\mathcal O_{X',x})=\dim(X')-1$ ($\mathbf 0$, 14.3.2), ce qui prouve que $S'$
est équicodimensionnel (5.1.4).
\end{env}

\begin{env}[10.7.3]
\label{IV.10.7.3-fr}
Soit $A$ un anneau de valuation discrète~; montrons que dans l'anneau
$B=A[T_1,\ldots,T_n]$ des polynômes à $n$ indéterminées, il existe deux idéaux
maximaux $\mathfrak m$, $\mathfrak n$ de \emph{hauteurs respectives} $n$ et
$n+1$. On l'a vu dans (5.2.5, (i)) pour $n=1$~; démontrons-le par récurrence sur
$n$. Comme $A[T_1,\ldots,T_{n+1}]$ est un $A[T_1,\ldots,T_n]$-module libre,
donc fidèlement plat, il y a dans $A[T_1,\ldots,T_{n+1}]$ deux idéaux maximaux
$\mathfrak m'$, $\mathfrak n'$ respectivement au-dessus de $\mathfrak m$ et
$\mathfrak n$ ($\mathbf 0_{\mathrm I}$, 6.5.1)~; en outre, d'après (5.5.3), ces
idéaux sont nécessairement de hauteurs respectives $n+1$ et $n+2$, d'où notre
assertion. Supposons $n\geq2$ dans ce qui suit. Soit
$\mathfrak J=\mathfrak m\cap\mathfrak n=\mathfrak m\mathfrak n$, et
$R=1+\mathfrak J$ qui est une partie multiplicative de $B$~; si l'on pose
$B'=R^{-1}B$, l'idéal $\mathfrak J'=R^{-1}\mathfrak J$ est contenu dans le
radical de $B'$ (Bourbaki, \emph{Alg. comm.}, chap.~III, \S~3,
n\textsuperscript{o}~5, prop.~12)~; on sait que $X'=\operatorname{Spec}(B')$
s'identifie en tant qu'espace topologique à un sous-espace de
$X=\operatorname{Spec}(B)$, et qu'aux points $x$ de $X'$, les anneaux locaux
$\mathcal O_{X,x}$ et $\mathcal O_{X',x}$ sont les mêmes (\textbf{I}, 1.6.2).
Considérons alors dans $X'$ l'ensemble fermé $Y'=V(\mathfrak J')$, et soit
$S=X'-Y'$~; on sait (10.5.7) que $S$ est un préschéma de Jacobson, évidemment
irréductible et noethérien~; en outre les anneaux locaux
$\mathcal O_{S,s}=\mathcal O_{X,s}$ sont universellement caténaires pour tout $s\in S$
en vertu de (5.6.3), puisque $A$ est universellement caténaire (5.6.4). Pourtant,
il y a deux points fermés $a$, $b$ de $S$ tels que $\mathcal O_{S,a}$ et
$\mathcal O_{S,b}$ \emph{n'aient pas même dimension}, autrement dit $S$ ne vérifie
pas la condition $3^\circ$ de (10.6.1). Pour le voir, considérons les deux idéaux
maximaux $\mathfrak m'=R^{-1}\mathfrak m$, $\mathfrak n'=R^{-1}\mathfrak n$ de
$B'$, qui sont de hauteurs $n$ et $n+1$ respectivement~; on a
$\mathfrak J'=\mathfrak m'\cap\mathfrak n'$, et $\mathfrak J'$ n'est donc
contenu dans aucun idéal premier de $B'$ distinct de $\mathfrak m'$ et
$\mathfrak n'$, qui sont par suite les seuls idéaux maximaux de $B'$. Soient
$a'$, $b'$ les seuls points fermés de $X'$, correspondant à $\mathfrak m'$ et
$\mathfrak n'$. Il existe dans $B'$ un idéal premier non maximal
$\mathfrak p'\subset\mathfrak m'$ qui n'est pas contenu dans $\mathfrak n'$~:
il suffit de montrer qu'il y a dans $B$ un idéal premier non maximal contenu dans
$\mathfrak m$ et non dans $\mathfrak n$~; pour cela, on pourra par exemple
considérer la fibre de $\mathfrak m$ pour le morphisme correspondant à l'injection
$A[T_1]\to B=A[T_1,\ldots,T_n]$, et appliquer (6.1.2). En considérant une chaîne
maximale d'idéaux premiers entre $\mathfrak p'$ et $\mathfrak m'$ et remplaçant
$\mathfrak p'$ par l'avant-dernier idéal de cette chaîne, on peut donc supposer
que le point $a$ de $X'$ correspondant à $\mathfrak p'$ est tel que son adhérence
dans $X'$ soit $\{a,a'\}$~; comme $B'_{\mathfrak m'}$ est biéquidimensionnel,
$\mathfrak p'$ est alors de hauteur $n-1$. On construit de la même manière un
idéal premier non maximal $\mathfrak q'$ de $B'$ de hauteur $n$, tel que si $b$
est le point correspondant de $X'$, l'adhérence de $b$ dans $X'$ soit
$\{b,b'\}$. Cela étant, $a$ et $b$ sont dans $S$, donc \emph{fermés dans $S$},
et répondent par suite à la question.
\end{env}

\subsection{Profondeur rectifiée}
\label{subsection:IV.10.8-fr}

\begin{definition}[10.8.1]
\label{IV.10.8.1-fr}
Soient $X$ un préschéma localement noethérien, $\mathcal F$ un $\mathcal O_X$-Module
cohérent. Pour tout $x\in X$, on appelle \emph{profondeur rectifiée de $\mathcal F$
au point $x$} et on note $\operatorname{prof}^{*}_{x}(\mathcal F)$ le nombre (entier
$\geq0$ ou $+\infty$) égal à
\begin{equation}
\label{IV.10.8.1.1-fr}
\tag{10.8.1.1}
\operatorname{prof}^{*}_{x}(\mathcal F)
=\operatorname{prof}(\mathcal F_x)+\dim(\overline{\{x\}})
\end{equation}
où $\overline{\{x\}}$ est l'adhérence du point $x$ dans $X$. Pour toute partie
$Z$ de $X$, on appelle \emph{profondeur rectifiée de $\mathcal F$ le long de $Z$}
et on note $\operatorname{prof}^{*}_{Z}(\mathcal F)$ le nombre
\begin{equation}
\label{IV.10.8.1.2-fr}
\tag{10.8.1.2}
\operatorname{prof}^{*}_{Z}(\mathcal F)
=\inf_{x\in Z}\operatorname{prof}^{*}_{x}(\mathcal F).
\end{equation}
\end{definition}

En d'autres termes pour tout entier $n$, la relation
$\operatorname{prof}^{*}_{Z}(\mathcal F)\geq n$ équivaut à
$\operatorname{prof}^{*}_{x}(\mathcal F)\geq n$ pour tout $x\in Z$. Si $Z=X$, on
écrit $\operatorname{prof}^{*}(\mathcal F)$ au lieu de
$\operatorname{prof}^{*}_{X}(\mathcal F)$.

